\chapter{Risk Controls}\label{ch:hf:risk-controls}

On 1 August 2012 a market maker's routing system sent millions of orders into the market in about forty-five minutes. It obtained over 4
million executions in 154 stocks, for more than 397 million shares, and lost more than \$460 million. The regulator's order afterwards reads as a
checklist of the controls it did not have. In this chapter's reconstruction, a \omterm{def:hf:risk-controls:loss}{loss limit} of \$2 million checked every second would have stopped
the same runaway after twelve seconds and \$2.1 million, without firing once on 2,500 ordinary days; an order-rate throttle, the control most
often named first, would not have stopped it at all.

\section{Pre-trade limits}

A market maker's orders pass through a pre-trade gate before they reach a venue: a set of checks that refuse any order that breaks a limit. One
Quant Book 13, chapter 22 builds the gate that runs these checks in nanoseconds; this chapter sets the policy it enforces.

\begin{definition}[Pre-trade risk check]\label{def:hf:risk-controls:pretrade}
A \emph{pre-trade risk check}\index{pre-trade risk check} is a test an order must pass before it is sent to a venue, against limits on its size,
value and price and on the positions, open orders and exposures it would create, refusing the order if any limit would be broken.
\end{definition}

\begin{definition}[Price collar]\label{def:hf:risk-controls:collar}
A \emph{price collar}\index{price collar} is a pre-trade limit on how far an order's price may be from a reference price (the last trade, the mid, a
fair value), in basis points or ticks, above which a buy or below which a sell is refused as a probable error.
\end{definition}

The limits form a hierarchy (\cref{fig:hf:risk-controls:hierarchy}): the firm, its desks, their strategies, the instruments. Each level caps what the
levels below it can add up to. The chapter's schema (\texttt{firm.riskctl}) writes it as a snapshot the gate can load: gross exposure and \omterm{def:hf:risk-controls:loss}{loss
limits} for the firm and each desk, an order rate, open orders, a \omterm{def:hf:risk-controls:loss}{loss limit} and a position limit for each strategy, and a collar, a maximum quantity
and notional and long and short limits for each instrument. A snapshot is validated before it is used (every desk within the firm's limits, every
strategy on a known desk) and exported with the gate's section, which Book 13's gate reads as it stands.

\begin{omfigure}
\begin{tikzpicture}[font=\footnotesize,>=stealth,box/.style={draw,rounded corners,fill=omThm!8,minimum width=2.4cm,minimum height=0.7cm,align=center}]
\node[box] (f) at (0,2.4) {firm\\gross, loss, messages};
\node[box] (d1) at (-2.2,1.2) {desk\\gross, loss};
\node[box] (d2) at (2.2,1.2) {desk\\gross, loss};
\node[box] (s1) at (-3.9,0) {strategy\\rate, open, loss};
\node[box] (s2) at (-0.9,0) {strategy};
\node[box] (s3) at (2.2,0) {strategy};
\node[box,fill=gray!10] (i) at (0,-1.2) {instruments: collar, size, notional, long, short};
\draw[->] (f) -- (d1); \draw[->] (f) -- (d2); \draw[->] (d1) -- (s1); \draw[->] (d1) -- (s2); \draw[->] (d2) -- (s3);
\node[box,fill=omDef!10] (g) at (6.8,1.2) {pre-trade gate\\(refuse)};
\node[box,fill=omDef!10] (m) at (6.8,0) {post-trade monitor\\(loss, position, rate)};
\node[box,fill=omDef!20] (k) at (6.8,-1.2) {kill switch\\(cancel, flatten, block)};
\draw[->] (m) -- (k);
\draw[->,dashed] (2.2,2.4) -- node[above] {limits snapshot} (6.8,2.4) -- (g);
\end{tikzpicture}

\omcaption{The limits hierarchy and the three places it acts: the pre-trade gate refuses an order that breaks a limit; the post-trade monitor watches
losses, positions and message rates and triggers the \omterm{def:hf:risk-controls:kill}{kill switch}, which cancels open orders, flattens positions and blocks new orders at the level
that broke. Source: \texttt{firm.riskctl}; the gate is One Quant Book 13's \texttt{firm.riskgate}.}\label{fig:hf:risk-controls:hierarchy}
\end{omfigure}

\section{Position, loss and message limits}

\begin{definition}[Loss limit]\label{def:hf:risk-controls:loss}
A \emph{loss limit}\index{loss limit} is a threshold on the realised plus marked-to-market loss of a strategy, desk or firm over a period (a day,
usually), beyond which the loss triggers a kill action at that level.
\end{definition}

Position and \omterm{def:hf:risk-controls:loss}{loss limits} are checked after trades as well as before them: the gate can refuse an order that would exceed a position, but only a
monitor that marks positions to the market sees a loss build (\cref{lst:hf:risk-controls:monitor}). A \omterm{def:hf:risk-controls:loss}{loss limit} is the most direct control of all,
since it limits what is being protected, but it acts only after the loss has happened and only as often as positions are marked. Message limits
(chapter 10's throttles and order-to-trade ratios) protect the venue and the firm's standing with it; the firm's own limit on messages a second is a
monitor on the whole firm, not a strategy.

\section{Price collars and fat-finger checks}

Collars and size checks catch the error in a single order: a price far from the market, a quantity or notional far above normal. On 2 May 2022 a
bank's trader who meant to sell a \$58 million basket created one of \$444 billion; the bank's controls blocked \$255 billion of it but let \$189
billion reach a trading algorithm, which sold about \$1.4 billion in European markets before it was cancelled. The UK's conduct regulator, fining
the bank £27.8 million in 2024, found that no hard block would have rejected the basket in its entirety, that the trader could close a pop-up alert
without reading it, and that real-time monitoring was too slow.

\begin{dated}{2026-09}{The rules on market access and algorithmic trading}\label{dat:hf:risk-controls:rules}
The SEC adopted Rule 15c3-5 in November 2010: brokers or dealers with market access must maintain risk management controls reasonably designed
to prevent erroneous orders and orders that exceed pre-set credit and capital thresholds, and their chief executive must certify them. Its order of
16 October 2013 fined the market maker of the 2012 runaway \$12 million for violating the rule. In the European Union, Commission Delegated
Regulation (EU) 2017/589 requires an investment firm to be able to cancel immediately, as an emergency measure, any or all of its unexecuted orders
on any or all venues, and to identify the algorithm and trader responsible for each order. The UK's Financial Conduct Authority fined a bank
£27,766,200 on 22 May 2024 over the 2022 basket error, and the Prudential Regulation Authority a further £33,880,000.
\end{dated}

\begin{definition}[Market access rule]\label{def:hf:risk-controls:access}
The \emph{market access rule}\index{market access rule} is SEC Rule 15c3-5, which requires a broker-dealer with access to an exchange or
alternative trading system, for itself or its customers, to have documented, regularly reviewed pre-trade controls against erroneous orders and
orders beyond pre-set credit and capital thresholds, under its direct and exclusive control.
\end{definition}

\section{Kill switches: who pulls them and when}

\begin{definition}[Kill switch]\label{def:hf:risk-controls:kill}
A \emph{kill switch}\index{kill switch} is a control that, when triggered by a person or by a monitor, stops a strategy, desk or the whole firm from
trading at once: it cancels its open orders, refuses new ones and, if so configured, flattens its positions, and it records who pulled it and why.
\end{definition}

The chapter's runaway (\cref{lst:hf:risk-controls:runaway}) takes the SEC order's figures: about 1,480 orders a second of 100 shares for 45
minutes, a loss of \$1.16 a share traded, and \$2.46 million a second of gross position (the \$3.5 billion long and \$3.15 billion short it
ended with). Stopping it flattens the position at 0.2\% of its value. Each control, alone, stops it at a different time
(\cref{fig:hf:risk-controls:controls}), and each fires on some share of ordinary days of a legitimate book (peak rates around 300 messages a
second, gross positions around \$150 million, intraday drawdowns around \$150,000, orders within tens of basis points of their references).

\begin{center}
\small
\begin{tabular}{@{}lrrr@{}}
\toprule
control & stops after & loss & fires on ordinary days\\
\midrule
none (the 45 minutes of 2012) & 2,700 s & \$477 million & --\\
a person at 5 minutes & 300 s & \$53 million & --\\
order-rate throttle, 500 a second & never & \$161 million & 10.2\%\\
\omterm{def:hf:risk-controls:collar}{price collar}, 5\% & 600 s & \$106 million & 0\%\\
capital threshold, \$1 billion gross & 407 s & \$72 million & 0\%\\
position limit, \$250 million gross & 102 s & \$18 million & 7.6\%\\
\omterm{def:hf:risk-controls:loss}{loss limit}, \$2 million, each second & 12 s & \$2.1 million & 0\%\\
all of them & 35 s & \$2.1 million & 10.2\%\\
\bottomrule
\end{tabular}
\end{center}

\begin{omfigure}
\begin{tikzpicture}
\begin{axis}[width=11cm,height=5.4cm,ybar,bar width=12pt,ymode=log,log basis y=10,ymin=1,ymax=900,xmin=-0.6,xmax=7.6,
  xtick={0,1,2,3,4,5,6,7},xticklabels={none,{person, 5 min},throttle,collar,capital,position,{loss limit},all},
  xticklabel style={font=\scriptsize,rotate=25,anchor=north east},ylabel={\small runaway's loss (\$ million)},
  tick label style={font=\footnotesize},grid=major,grid style={gray!25},table/col sep=comma,nodes near coords,
  nodes near coords style={font=\scriptsize},point meta=rawy]
\addplot[fill=omThm!60,draw=omThm] table[x=i,y=loss]{figdata/hft/27-risk-controls/controls.csv};
\end{axis}
\end{tikzpicture}

\omcaption{The loss of a runaway shaped on the 2012 incident under each control alone and under all of them together (log scale, \$ million):
no control (stopped by hand at 45 minutes), a person at 5 minutes, an order-rate throttle of 500 a second, a 5\% \omterm{def:hf:risk-controls:collar}{price collar}, a \$1 billion capital
threshold, a \$250 million position limit and a \$2 million \omterm{def:hf:risk-controls:loss}{loss limit} marked each second. Data: \texttt{hf\_risk.table}.}\label{fig:hf:risk-controls:controls}
\end{omfigure}

The throttle, often the first control named, does not stop anything: it slows the runaway to a third and the loss follows, \$161 million by the end.
The \omterm{def:hf:risk-controls:collar}{price collar} and the capital threshold, the controls the \omterm{def:hf:risk-controls:access}{market access rule} names, stop it after ten and seven minutes at \$106 and \$72
million. The position limit stops it in under two minutes; the \omterm{def:hf:risk-controls:loss}{loss limit}, in twelve seconds. Together they stop it at 35 seconds: the throttle,
by slowing the runaway, delays the \omterm{def:hf:risk-controls:loss}{loss limit}, but the loss is the same.

Each threshold is a trade-off (\cref{fig:hf:risk-controls:tradeoff}): a tighter \omterm{def:hf:risk-controls:loss}{loss limit} stops the runaway sooner but fires on ordinary days, 2\%
of them at \$1 million, 19\% at \$500,000; a position limit tight enough to stop it in a minute fires on half of all days. The \omterm{def:hf:risk-controls:loss}{loss limit} dominates:
for the same false alarms it stops the runaway much sooner, because it measures what matters. The position limit catches what the \omterm{def:hf:risk-controls:loss}{loss limit}
cannot: a runaway that has not yet lost.

\begin{omfigure}
\begin{tikzpicture}
\begin{axis}[width=11cm,height=5.4cm,xlabel={\small ordinary days on which the control fires (\%)},ylabel={\small runaway's loss (\$ million)},
  xmin=-2,xmax=92,ymode=log,log basis y=10,ymin=0.4,ymax=60,tick label style={font=\footnotesize},grid=major,grid style={gray!25},
  table/col sep=comma,legend style={font=\scriptsize,at={(0.97,0.97)},anchor=north east}]
\addplot[omThm,thick,mark=*,mark size=1.4pt] table[x=false,y=loss]{figdata/hft/27-risk-controls/sweep_loss.csv};
\addplot[omDef,thick,dashed,mark=square*,mark size=1.3pt] table[x=false,y=loss]{figdata/hft/27-risk-controls/sweep_position.csv};
\legend{{loss limit, \$0.5--4 million},{position limit, \$100--600 million}}
\end{axis}
\end{tikzpicture}

\omcaption{The trade-off of a threshold: the runaway's loss under a \omterm{def:hf:risk-controls:loss}{loss limit} (\$0.5 to \$4 million) and under a gross position limit (\$100 to
\$600 million) against the share of 2,500 ordinary days on which each would have fired. Data: \texttt{hf\_risk.sweep}.}\label{fig:hf:risk-controls:tradeoff}
\end{omfigure}

Who pulls the switch matters as much as when. A monitor fires in a second; a person takes minutes to see, understand and act, and in 2012 took
forty-five. The \omterm{def:hf:risk-controls:kill}{kill switch} must be reachable by the monitor and by a person, must cancel before it flattens, must act on every venue at once (as
the EU's rule requires), and must leave an audit trail of who pulled it and why (\texttt{firm.riskctl.KillSwitch}).

\section{The regulatory requirements for algorithmic firms}

The rules in the dated box converge on the same controls: pre-trade limits on price, size and value; limits on credit and capital; a kill
function; knowing which algorithm and trader sent each order; and review of the controls by management. The regulators' orders read as descriptions
of what went wrong when one was missing: in 2012, no control on the aggregate capital the orders committed, and no automated alert that led to action
in time; in 2022, a size check that could be overridden and monitoring too slow to act.

\section{Tutorial: forty-five minutes}\label{tut:hf:risk-controls:tut}

\begin{tutorial}
\textbf{Goal.} Write the limits snapshot, export it for the gate, and measure what each control would have saved against a 2012-shaped runaway and
what it costs on ordinary days. \textbf{End state:} the table and the three figures.
\begin{tutsteps}
\item \textbf{The snapshot}: \texttt{firm.riskctl.validate} and \texttt{export}; \texttt{riskctl\_fixture.py} writes
\texttt{data/limits.json}, which One Quant Book 13's \texttt{firm.riskgate.from\_riskctl} loads.
\item \textbf{The monitor}: marks, \omterm{def:hf:risk-controls:loss}{loss limits} and kill actions at each level.
\omcode{code/firm/riskctl/firm_riskctl.py}{126}{143}{Marks positions to the market, sums strategies into desks and the firm, and kills at the level whose loss limit breaks.}\label{lst:hf:risk-controls:monitor}
\item \textbf{The runaway} under each control.
\omcode{code/firm/riskctl/firm_riskctl.py}{169}{188}{Each control's stopping time; the first one stops the runaway, which is then flattened.}\label{lst:hf:risk-controls:runaway}
\item \textbf{Ordinary days} and the threshold sweeps (\texttt{hf\_risk.table}, \texttt{hf\_risk.sweep}).
\end{tutsteps}
\textbf{What to change next.} Replay the runaway through Book 13's gate on \texttt{firm.exchsim} (its \texttt{firm\_riskgate\_runaway}); add a
second runaway that loses slowly and see which control catches it; make the \omterm{def:hf:risk-controls:loss}{loss limit} intraday-drawdown based.
\end{tutorial}

\section{Build: the risk controls}\label{bld:hf:risk-controls:riskctl}

\begin{build}
\bfield{Purpose} Own the limits policy (hierarchy, schema, validation), the post-trade monitors and the \omterm{def:hf:risk-controls:kill}{kill switch}; hand the pre-trade gate its
snapshot.
\bfield{Interface} \texttt{validate(limits)}, \texttt{to\_riskgate(limits)}, \texttt{export(limits)}, \texttt{KillSwitch}, \texttt{Monitor(limits,
switch)} with \texttt{on\_fill}, \texttt{mark}, \texttt{on\_message}; \texttt{runaway(controls)}, \texttt{normal\_days(controls)}. The schema,
version 1, is in the module's docstring and is the contract with Book 13's \texttt{firm.riskgate}.
\bfield{Rules} Dollars, shares, nanoseconds; a desk within the firm's limits; kills recorded with their reason; cancel before flatten.
\bfield{Acceptance tests} \texttt{code/firm/riskctl/tests/}: validation catches a desk above the firm and an unknown desk; the gate section by hand;
the fixture equals the export and Book 13's gate loads it and refuses an order outside the collar; the monitor kills a strategy beyond its \omterm{def:hf:risk-controls:loss}{loss limit}
and the firm beyond its message rate, and the audit records kill and unkill; the runaway and ordinary days behave as in the chapter.
\bfield{Stretch} Intraday drawdown limits; per-venue kill; limits changed during the day with versioned snapshots.
\end{build}

\begin{omsources}
\item US Securities and Exchange Commission, \emph{In the Matter of Knight Capital Americas LLC}, Release No.\ 34-70694, 16 October 2013.
\item Commission Delegated Regulation (EU) 2017/589, Article 12 (kill functionality).
\item Financial Conduct Authority, FCA fines CGML £27,766,200 for failures in its trading systems and controls, 22 May 2024.
\end{omsources}

\section{Exercises}

\begin{exercise}[$\star$]\label{exo:hf:risk-controls:1}
The runaway loses \$1.16 a share on 148,100 shares a second. How long does a \$2 million \omterm{def:hf:risk-controls:loss}{loss limit}, checked every second, take to fire?
\end{exercise}

\begin{exercise}[$\star$]\label{exo:hf:risk-controls:2}
Gross position grows by \$2.46 million a second. When does a \$250 million position limit stop it, and a \$1 billion capital threshold?
\end{exercise}

\begin{exercise}[$\star$]\label{exo:hf:risk-controls:3}
A buy order is priced at \$42.10 against a reference of \$40.00. Does a 5\% collar refuse it?
\end{exercise}

\begin{exercise}[$\star\star$]\label{exo:hf:risk-controls:4}
Why does the throttle not stop the runaway, and what is it for?
\end{exercise}

\begin{exercise}[$\star\star$]\label{exo:hf:risk-controls:5}
Why does adding every control stop the runaway later (35 seconds) than the \omterm{def:hf:risk-controls:loss}{loss limit} alone (12)?
\end{exercise}

\begin{exercise}[$\star\star$]\label{exo:hf:risk-controls:6}
What should a \omterm{def:hf:risk-controls:kill}{kill switch} do first: cancel or flatten? Why?
\end{exercise}

\begin{exercise}[$\star\star\star$]\label{exo:hf:risk-controls:7}
\emph{Coding.} With \texttt{hf\_risk.sweep}, which \omterm{def:hf:risk-controls:loss}{loss limit} fires on no more than 1\% of ordinary days, and what does it let the runaway lose?
\end{exercise}

\begin{exercise}[$\star\star\star$]\label{exo:hf:risk-controls:8}
\emph{Find the flaw.} ``Our algorithm has never lost more than \$500,000 in a day, so a \$500,000 \omterm{def:hf:risk-controls:loss}{loss limit} costs us nothing.''
\end{exercise}

\section{Problem: Forty-Five Minutes}

\begin{problem}[{Weekend problem --- forty-five minutes}]\label{pb:hf:risk-controls:1}
A market-making firm sets its risk controls after reading the 2012 order.

\medskip\noindent\textbf{Part I --- The incident and the rules.}
\begin{enumerate}
\item Summarise what happened on 1 August 2012, by the SEC's order.
\item Define the \omterm{def:hf:risk-controls:access}{market access rule} and what it requires.
\item What does the EU's kill functionality require?
\item What went wrong in the 2022 basket error, by the FCA's account?
\end{enumerate}

\medskip\noindent\textbf{Part II --- The controls.}
\begin{enumerate}[resume]
\item Define a \omterm{def:hf:risk-controls:pretrade}{pre-trade risk check} and a \omterm{def:hf:risk-controls:collar}{price collar}.
\item Define a \omterm{def:hf:risk-controls:loss}{loss limit} and a \omterm{def:hf:risk-controls:kill}{kill switch}.
\item Describe the limits hierarchy and the snapshot's validation.
\item How does the snapshot reach the nanosecond gate?
\end{enumerate}

\medskip\noindent\textbf{Part III --- Measurements.}
\begin{enumerate}[resume]
\item Give the stopping time and loss under each control.
\item Which controls fire on ordinary days, and how often?
\item Describe the loss-limit and position-limit trade-offs.
\item Why does the \omterm{def:hf:risk-controls:loss}{loss limit} dominate, and what does the position limit add?
\end{enumerate}

\medskip\noindent\textbf{Part IV --- The verdict.}
\begin{enumerate}[resume]
\item State the \emph{named result}: the loss the runaway would have made under each control, and the earliest control that stops it.
\item Which set of controls would you run, with which thresholds?
\item Who should be able to pull the \omterm{def:hf:risk-controls:kill}{kill switch}, and how fast?
\item What must the audit record?
\item What does a runaway that loses slowly need?
\item How would you test the controls without a runaway?
\item What should a firm's chief executive certify?
\item In one sentence: what does a risk control buy?
\end{enumerate}
\end{problem}

\section{Interview questions}

\begin{interviewq}[$\star$ \iqroles{risk}]\label{iq:hf:risk-controls:1}
List the pre-trade checks you would require on every order, in the order you would run them.
\end{interviewq}

\begin{interviewq}[$\star\star$ \iqroles{developer}]\label{iq:hf:risk-controls:2}
Your gate refuses orders when its limits snapshot is stale. Why fail closed, and what could go wrong?
\end{interviewq}

\begin{interviewq}[$\star\star$ \iqroles{trader}]\label{iq:hf:risk-controls:3}
Your strategy hit its \omterm{def:hf:risk-controls:loss}{loss limit} at 10:05 on a volatile day and was flattened. The market recovered by 10:30. Was the limit wrong?
\end{interviewq}

\begin{interviewq}[$\star\star$ \iqroles{risk}]\label{iq:hf:risk-controls:4}
How would you set a strategy's position limit from its history?
\end{interviewq}

\begin{interviewq}[$\star\star$ \iqroles{developer}]\label{iq:hf:risk-controls:5}
Design the \omterm{def:hf:risk-controls:kill}{kill switch} for a firm connected to 20 venues through 3 gateways. What must happen in the first millisecond?
\end{interviewq}

\begin{interviewq}[$\star\star\star$ \iqroles{researcher}]\label{iq:hf:risk-controls:6}
With ordinary daily drawdowns normal with mean $\mu$ and standard deviation $\sigma$, and a runaway losing at rate $\ell$ per second, find the \omterm{def:hf:risk-controls:loss}{loss
limit} that minimises the expected cost of false alarms plus the runaway's expected loss, given a runaway once in $N$ days.
\end{interviewq}
